% Đoạn LaTeX có thể chèn trực tiếp vào báo cáo kỹ thuật AIC 2026.
% Yêu cầu preamble gợi ý: \usepackage[utf8]{inputenc}, \usepackage{booktabs}.

\subsection{Truy vấn hỏi--đáp đa phương thức}

Nhánh hỏi--đáp (Q\&A) được thiết kế theo hai tầng độc lập: truy xuất bằng chứng
và đọc câu trả lời. Trước hết, bộ phân tích luật nhận diện ý định của câu hỏi
(chẳng hạn đếm đối tượng, đọc văn bản, nhận diện người hoặc địa điểm) và kiểu
đáp án mong đợi. Hệ thống sau đó hợp nhất tín hiệu CLIP-L, OCR, ASR và caption
để xếp hạng các khung hình ứng viên; ngân sách truy xuất được điều chỉnh theo
độ khó của truy vấn. Mỗi ứng viên luôn giữ cặp \texttt{video\_id} và
\texttt{frame\_idx} thật từ bảng ánh xạ keyframe, tránh nhầm số thứ tự ảnh
\texttt{n} với chỉ số khung hình cần nộp.

Ở tầng đọc, hệ thống hỗ trợ hai chế độ. Chế độ cục bộ khai thác OCR, ASR và bộ
đọc ảnh nhẹ. Chế độ đám mây tùy chọn gửi tối đa một tập nhỏ ảnh ứng viên cùng
ngữ cảnh OCR/ASR/caption đến Gemini để ước lượng độ liên quan và sinh câu trả
lời ngắn cho từng ảnh. Kết quả được ép theo JSON schema, kiểm tra ánh xạ ngược
về đúng ứng viên, sửa lỗi JSON có kiểm soát và thử lại tối đa một lần. Phản hồi
được cache theo câu hỏi, model, phiên bản prompt và định danh bằng chứng. Khi
dịch vụ không khả dụng, pipeline trả về nhánh cục bộ và gắn cờ fallback thay vì
dừng toàn bộ giao diện. API key chỉ được đọc từ biến môi trường và không được
ghi vào mã nguồn, manifest hay báo cáo.

\subsubsection{Đánh giá và chống rò rỉ dữ liệu}

Chúng tôi tách hai chế độ đánh giá. Chế độ \texttt{r4} đo toàn tuyến từ danh
sách ứng viên QA-R4 đến công thức R@K của Ban Tổ chức; ground truth chỉ được mở
sau khi đã sinh tệp submission. Chế độ \texttt{oracle\_gt} chỉ dùng vị trí
video và khoảng frame chuẩn để chọn ảnh, nhưng không truyền câu trả lời chuẩn
vào prompt hoặc reader. Chế độ thứ hai chỉ dùng chẩn đoán năng lực đọc ảnh và
luôn được đánh dấu \texttt{diagnostic\_only}; nó không phải điểm end-to-end.
Các request bị lỗi API hoặc chuyển sang local fallback không được đưa vào điểm
trung bình của cấu hình Gemini. Tương tự, truy vấn không có keyframe trích sẵn
nằm trong khoảng GT được ghi nhận là \texttt{oracle\_keyframe\_gap} và loại
khỏi mẫu đánh giá reader.

\subsubsection{Kết quả chẩn đoán hiện tại}

Trên 18 truy vấn Q\&A của tập phát triển, artifact QA-R4 hiện tại đạt video
recall $5/18$ (27.78\%) và frame recall $1/18$ (5.56\%), cho thấy nút thắt
chính vẫn nằm ở truy xuất bằng chứng. Trong thí nghiệm oracle trên tập tune,
10/12 truy vấn có ít nhất một keyframe nằm trong khoảng GT; hai truy vấn còn
lại được phân loại là lỗi coverage. Một phép thử đơn lẻ trên truy vấn 44 cho
điểm 1.0 khi Gemini nhận đúng khung hình. Ở bốn truy vấn oracle được xử lý
thành công trước khi chạm quota dịch vụ, cấu hình 4 ảnh và 12 ảnh cùng đạt
3/4 câu đúng, trong khi cấu hình 8 ảnh đạt 1/4. Do kích thước mẫu đám mây còn
nhỏ, các con số này chỉ là kết quả chẩn đoán sơ bộ và chưa được dùng để khẳng
định cấu hình cuối cùng.

Kết quả trên dẫn đến quyết định kỹ thuật an toàn: giữ reader đám mây như một
thành phần tùy chọn có kiểm soát, không phụ thuộc vào nó để đảm bảo pipeline
hoạt động, và ưu tiên cải thiện video/frame retrieval trước khi mở rộng số lần
gọi API. Việc lựa chọn 4 hay 12 ảnh chỉ được khóa sau khi toàn bộ mẫu hợp lệ
được đánh giá mà không có fallback.
